\section{Evaluation}
HCMO is evaluated at four complementary levels: common ontology pitfalls and FAIR metadata, OWL consistency, instance-data validation, and executable competency questions. Reports and raw outputs are retained in the repository so that historical scores are not silently attributed to a later artifact.

\textbf{Pitfalls and FAIRness.} Historical scans of the clean-v2 precursor are kept as remediation evidence but are not attributed to the current graph. On 15 August 2026, OOPS! was rerun directly on the candidate RDF/XML distribution. It reported P04, P08, P10, P13, P22, P34, and P35. The important P10 finding is the scanner's ontology-wide missing-disjointness heuristic and names no affected pair. P34 and P35 concern reused SOSA, SemTS, QUDT, Schema.org, and PROV terms whose authoritative declarations are covered by HCMO's pinned external-vocabulary contract rather than copied into the merged graph. The minor P04, P08, and P22 findings likewise concern reused BFO or external terms. P13 marks 32 properties as lacking inverse declarations; candidate inverses are not asserted without definition-level evidence because doing so would create false entailments merely to satisfy a scanner.

FOOPS! was rerun by uploading the canonical candidate RDF/XML distribution and scored 1.0. All ontology metadata, licensing, provenance, resolvability, vocabulary-reuse, label, and definition tests passed. A separate deployment check found that the public PURL still served the older 1,252-triple graph rather than the 1,397-triple candidate, so the perfect candidate-file score is not misreported as a score for the currently deployed representation. Raw outputs and term-level triage are archived in \nolinkurl{docs/paper/evaluation/CANDIDATE-OOPS-FOOPS-2026-08-15.md}. Both scanners will be rerun after the candidate receives its final version IRI, tag, and PURL target.

\textbf{Logical consistency.} The canonical RDF/XML artifact was rebuilt and checksummed on 15 August 2026. \nolinkurl{dist/hcmo.owl} contains 1,397 RDF triples; the generated profile contains 60 declared release classes, 68 object properties, and 75 datatype properties, including compatibility and directly referenced external terms. HermiT loaded 65 classes and reported zero inconsistent classes. The pre-check also found no active \nolinkurl{UNKNOWN:} IRI and no property declared as both object and datatype. The exact checksum, environment, output, and triage are archived in \nolinkurl{docs/paper/evaluation/POST-PR24-26-REVIEW-2026-08-15.md}.

\textbf{SHACL validation.} The validator supplies the canonical ontology as a separate ontology graph and enables RDFS entailment when validating each example. This allows domain, range, and subclass consequences to select the same \nolinkurl{sh:targetClass} constraints as explicit types. Three positive examples conform, while two edge-case examples produce their reviewed violations. One negative fixture deliberately omits explicit \nolinkurl{rdf:type}: it conforms when validated without the ontology, then fails when ontology-aware inference selects the intended targets. This differential test guards the validator's entailment contract while keeping profile requiredness in SHACL rather than translating it into OWL existential semantics. A separate ISA/STATO evidence shape validates the pinned workflow boundary; an injected process/data cycle is required to fail.

The accepted round-trip fixture contains a 2 $\times$ 2 treatment-by-enrichment design with eight individually housed animals, 56 repeated dark-phase activity observations, non-overlapping re-housing, an actual derived tissue Sample, and separate fitted-model, estimate, confidence-interval, p-value, File, and file-fragment entities. Canonical HCMO RDF and extended ISA RO-Crate JSON-LD are graph-isomorphic at 1,539 triples. Dark-phase outcomes use phenomenon intervals, and the re-housed animal's observations resolve to its new enclosure through the authoritative assignment. The declared mixed model is executed against the generated CSV with pinned numerical dependencies; the model serialization and the active-versus-vehicle contrast at standard enrichment use distinct exact row fragments. Dedicated shapes and injected probes reject overlapping, orphaned, reversed, or incompletely generated housing assignments; generated replacement animals; observation/ housing mismatch; group/factor disagreement; and fabricated ISA assignment results.

\textbf{Competency questions.} Eleven SPARQL queries run over the ontology plus all positive examples. Negative fixtures are isolated from the query graph. The validator checks both query-index completeness and the complete expected answer rows, including bindings, unbound values, and multiplicity:

\begin{table}[t]
\centering
\small
\begin{tabularx}{\linewidth}{@{}Xr@{}}
\toprule
Competency question & Exact answers \\
\midrule
subjects assigned by monitored enclosure & 3 \\
ISA factor/group assignment & 1 \\
ISA raw-file recording provenance & 1 \\
OBI transformation and STATO result provenance & 1 \\
enclosures missing dimension records & 0 \\
environmental specification and observation quantities & 1 \\
current housing at the reviewed reference time & 1 \\
operational and calibration evidence at the reviewed time & 1 \\
housed subjects needing provisioning & 3 \\
properties captured by enclosure sensors & 4 \\
systems with observations lasting at least 24 hours under a condition & 1 \\
\bottomrule
\end{tabularx}
\end{table}

The zero-row dimensions query is intentional: it verifies absence of missing records in the positive fixture rather than being accepted as an unexamined empty result. The three ISA/STATO questions establish a narrow executable evidence slice; they do not establish HCMO class mappings or formal ISA RO-Crate conformance.

Five additional exact-answer questions run against the isolated round-trip fixture: nine housing-history rows, eight authoritative subject/group/factor rows, eight per-animal repeated-observation counts, four statistical-result/ file-fragment rows including exact values and meaning links, and one Source-to-derived-Sample row. \nolinkurl{roc-validator} 0.11.3 passes all RO-Crate 1.2 required rules and all ISA-specific required rules. Its full inherited ISA run has one isolated upstream conflict: the embedded ISA profile requires RO-Crate 1.1 although its minimal ISA fixture uses the 1.2 context. We report this as interoperability evidence, not formal conformance. Separately, pinned \nolinkurl{isatools} 0.14.3 validates an executable native projection and round trip for the genuine Source-to-Sample collection. Exact HCMO identifiers survive as ISA comments. The test deliberately excludes non-native housing, direct whole-animal assay, source-bound factor assignment, group, repeated-observation, and semantic-result/file-fragment structures and checks their loss manifests.

\begin{figure}[t]
\centering
\begin{tikzpicture}[
  node distance=5mm and 5mm,
  box/.style={draw, rounded corners=1.5pt, align=center, font=\scriptsize, inner sep=4pt},
  flow/.style={-{Latex[length=1.8mm]}, font=\tiny}
]
\node[box, fill=red!8] (animal) {animal Source};
\node[box, fill=yellow!12, right=of animal] (housing) {housing history\\9 reviewed rows};
\node[box, fill=blue!8, below=of animal] (obs) {56 repeated\\observations};
\node[box, fill=violet!8, right=of obs] (model) {fitted model\\STATO results};
\node[box, fill=green!9, right=of model] (files) {Files and exact\\row fragments};
\node[box, fill=red!8, below=of files] (sample) {derived tissue Sample};
\draw[flow] (animal) -- (housing);
\draw[flow] (animal) -- (obs);
\draw[flow] (obs) -- (model);
\draw[flow] (model) -- (files);
\draw[flow] (animal) to[bend right=20] node[below]{collection} (sample);
\end{tikzpicture}
\caption{Executable 2 $\times$ 2 fixture used to test identity, housing,
repeated observations, statistical semantics, and ISA projection boundaries.}
\label{fig:roundtrip}
\end{figure}

